\documentclass[12pt]{article}

\usepackage[T1]{fontenc}
\usepackage[utf8]{inputenc}
\usepackage[romanian]{babel}

\usepackage{amsmath, amssymb, amsthm}
\usepackage{mathrsfs} % pentru S caligrafic
\usepackage{geometry}
\geometry{a4paper, margin=2.5cm}

\title{Distribuția unei Variabile Aleatoare și Echivalența în Distribuție}
\author{}
\date{}

\theoremstyle{plain}
\newtheorem{definition}{Definiție}
\newtheorem{proposition}{Propoziție}
\newtheorem{theorem}{Teoremă}
\newtheorem{corollary}{Corolar}

\begin{document}

\maketitle

\section{Distribuția unei variabile aleatoare}

\begin{definition}
Fie $(\Omega, \mathcal{F}, \mathbb{P})$ un spațiu de probabilitate și
$X : \Omega \to S$ o variabilă aleatoare cu valori într-un spațiu măsurabil
$(S, \mathscr{S})$.
\textbf{Distribuția} (sau \textbf{legea}) lui $X$ este măsura de probabilitate
$\mathbb{P}_X$ definită pe $(S, \mathscr{S})$ prin:


\[
\mathbb{P}_X(A) := \mathbb{P}(X \in A),
\qquad A \in \mathscr{S}.
\]


\end{definition}

Distribuția lui $X$ este imaginea măsurii $\mathbb{P}$ prin aplicația $X$:


\[
\mathbb{P}_X = \mathbb{P} \circ X^{-1}.
\]



\begin{proposition}
Distribuția $\mathbb{P}_X$ este o măsură de probabilitate pe
$(S, \mathscr{S})$ și este determinată în mod unic de valorile sale pe
o familie generatoare a lui $\mathscr{S}$.
\end{proposition}

\section{Echivalența în distribuție}

\begin{definition}
Două variabile aleatoare $X, Y : \Omega \to S$ se numesc
\textbf{egale în distribuție} sau \textbf{echivalente în distribuție} dacă
au aceeași distribuție, adică:


\[
\mathbb{P}_X = \mathbb{P}_Y.
\]


Scriem:


\[
X \overset{d}{=} Y.
\]


\end{definition}

\begin{proposition}
Pentru două variabile aleatoare $X, Y : \Omega \to S$, următoarele afirmații sunt echivalente:
\begin{enumerate}
    \item $X \overset{d}{=} Y$;
    \item $\mathbb{P}(X \in A) = \mathbb{P}(Y \in A)$ pentru orice $A \in \mathscr{S}$;
    \item pentru orice funcție măsurabilă $g : S \to \mathbb{R}$ avem:
    

\[
    \mathbb{P}(g(X) \le t) = \mathbb{P}(g(Y) \le t)
    \quad \text{pentru toate } t \in \mathbb{R}.
    \]


\end{enumerate}
\end{proposition}

\begin{corollary}
Dacă $X \overset{d}{=} Y$, atunci pentru orice funcție măsurabilă
$g : S \to T$ între două spații măsurabile $(S, \mathscr{S})$ și
$(T, \mathcal{T})$ avem:


\[
g(X) \overset{d}{=} g(Y).
\]


\end{corollary}

\end{document}
